%% % % % % % % % % % % % % % % % % % % % % % % % % % % % % % %
%% % % % % % % % % % % % % % % % % % % % % % % % % % % % % % %
\documentclass[a4,11pt]{cip-v3-cls-submit}
\include{cip-v3-pack-and-defined}
\begin{document}

%PAGE_NUMBER%%%%%%%%%%%%%%
\Year{2023}
%\Volume{??}\Number{??}	\ArticleId{????}\Year{2023}\Page{1}\Endpage{10}
%%%%%%%%%%%%%%%%%%%%%%%%%%%%%%%%%%%%%%%%%%%%%
	\title{Guidelines for preparing manuscript for ``Communications in Physics"}
	%%%%%%%%%%%%%%%%%%%%%%%%%%%%%%%%%%%%%%%%%%%%%%%%%
	%%%%%AUTHORS AND ADDRESSES
	%%%%%%%%%%%%%%%%%%%%%%%%%%%%%%%%%%%%%%%%%%%%%%%%%
	\authors{
	\authord{A. B. Last-Name}{1,2}, \authord{Full Name}{1}, \authord{First Corresponding Author}{1,\dagger}
	\newline
\and \authord{Second Corresponding Author}{2,\ddagger}
	%%%%%%%%%%%%%%%%%%%
	\newline
	%%%%%%%%%%%%%%
	\addressed{1}{Institute of Physics, Vietnam Academy of Science and Technology
		10 Dao Tan, Hanoi 11108, Vietnam}
	\addressed{2}{Vietnam National University, Hanoi, 144 Xuan Thuy, Cau Giay, Hanoi 123105, Vietnam}
%\CorrEmail{$^{\dagger}$cip@vjs.ac.vn}
%\CorrEmail{$^{\ddagger}$abc@iop.vast.vn}
%\CorrEmail{cip@vjs.ac.vn}%\dagger
\CorrEmailDouble{cip@vjs.ac.vn}{abc@iop.vast.vn}%\dagger\ddagger
%\CorrEmailTriple{first@doman.com}{second@doman.com}{third@doman.com}%\dagger\ddagger\star
%%%%%%%%%%%%%%%%%%%%%
\PublicationInformation{\today}{(Editor will insert the date when the article is accepted)}{(Editor will insert the date when the article is published)}
%%%%%%%%%%%%%%
}
\maketitle
%%%%%%%%%%%%%%%%%%%%%%%%%%%%%%%%%%%%%%%%%%%%%%%%%%%%%%%%%%%%%%%%%%%%%%
	\markboth{Title of the article $\ldots$}{First Author \textit{et al.}}
%%%%%%%%%%%%%%%%%%%%%%%%%%%%%%%%%%%%%%%%%%%%%%%%%%%%%%%%%%%%%%%%%%%%%%


\begin{abstract}
This guide represents a LaTeX template, which helps authors when preparing a manuscript for  "Communications in Physics". This LaTeX template has been arranged so it can also be used as a reference manual for the manuscript preparation by Microsoft Word. Authors are encouraged to download the LaTex template at \href{https://vjs.ac.vn/index.php/cip/libraryFiles/downloadPublic/16}{\textbf{\textcolor{red}{CIP-V3-Submit-Template}}} and refer for manuscript preparation before submission to ``Communications in Physics". This LaTeX template can be compiled by PdfLaTex or standard LaTeX. 
\end{abstract}

\Keywords{Communications in Physics, \LaTeX, template}

\Classification{classification categories, Physics and Astronomy Classification Scheme}


\section{Introduction}

The article manuscript should be prepared in English and generally be arranged in the  following order:  title  page, main body and References.  The title page consists of Title, Authors names, Authors affiliations, Receipt Date, Abstract, Keywords, and Classification. The main body of the manuscript consists of Introduction, Main sections, Conclusions, Acknowledgements.
When necessary the statement of the authors contributions and the statement of conflict of interest can be placed after Acknowledgements, and before References. Appendices if any should be placed after References.


\section{Page format}

All text must be in a one-column format.

The length of a research article can vary from 6 to 15 pages.
The length of a review article can vary from 6 to 20 pages
The manuscript with the length beyond these page limits will not be reviewed.


If the manuscript is prepared by Microsoft Word, the Times New Roman font, size 12pt and the back color should be used. The page size is 19.5 cm $\times$ 26.6 cm. The top and bottom margins are 2.86 cm, and the left and right margins are 2.86 cm.



\section{Title page}

\subsection{Title}
Title is to be  descriptive and comprehensible. Begin the first word with a capital letter, and thereafter capitalize only proper or trade names. The mathematical, physical, chemical symbols if any should be written in the standard scientific notations.

If the manuscript is prepared by Microsoft Word, the Title is written by the Times New Roman, size 14 pt and bold face.

\subsection{Authors names}
Begin with initials of the first and middle name with period and space followed by the surname, for example, L. D. Landau , P. Fulde, X. V. Tran. The authors names can also be written in the full common form, for example, Peter Fulde, Dam T Son, Tran Van Xuan. It is preferable to use only one form of the authors names in all publications. Clearly indicate the corresponding authors who will handle correspondence at all stages of refereeing, publication, and post-publication by email addresses.

If the manuscript is prepared by Microsoft Word, the authors names are written by Times New Roman,  size 12 pt and bold face.

\subsection{Authors affiliations}
The affiliations are the places at which the authors conducted the research that they have reported in the manuscript. The affiliations should follow this order: department/laboratory/section/ division, institution, city with postal code (if any), country. The street numbers and names are written only if necessary for effective mail delivery.
Use superscripts to match authors with institutions.

\subsection{Receipt date}
The receipt date indicates the date the manuscript was received by the scientific editor. This date will appear in the printed article.

\subsection{Abstract}
The abstract should consist of one paragraph and be completely self-contained.
Non-standard or uncommon abbreviations or references must be avoided.

\subsection{Keywords}
Include at least three keywords, that best describe the article content. They are in lower case, except for proper or trade names, and in alphabetical order, separated by commas.

\subsection{Classification}
Select the classification categories according to Physics and Astronomy Classification Scheme, which is
available at https://vjs.ac.vn/index.php/cip/libraryFiles/downloadPublic/18.


\section{Main body}
The main body of the manuscript consists of Introduction, main sections, Conclusions, Acknowledgements,
and when necessary it also includes the statement of the authors contributions and the statement of conflict of interest.

In the manuscript all references, figures, and tables must be cited consecutively in numerical order.
References use arabic numerals.
Figures and tables use arabic numerals with (a), (b), (c), etc. to label their parts, for example,
figure 1(a), table 2(b), etc.


\subsection{Formulas}
Mathematical, physical, chemical symbols and formulas should clearly be typed in the standard scientific notations.
Sort mathematical, physical, chemical ... formulas can be written within the text lines, while
long ones can be written separately in new lines. The sort formulas are not numbered. The long formulas
are numbered consecutively with arabic numbers within parentheses on the right-hand side of each formula. Only cited long formulas are numbered. The long formulas can grammatically have dots or commas at their ends like in the correct sentences grammar. For example, single equation looks like
\begin{equation}
\mathbf{G}(\mathbf{k},\omega)=\mathbf{G}_0(\mathbf{k},\omega)+
\mathbf{G}_0(\mathbf{k},\omega) \boldsymbol{\Sigma}(\mathbf{k},\omega) \mathbf{G}(\mathbf{k},\omega) .
\end{equation}
The system of equations looks like
\begin{eqnarray}
\mathbf{G}_0(\mathbf{k},\omega) &=& \frac{1}{\omega-\varepsilon(\mathbf{k})} , \\
\boldsymbol{\Sigma}(\mathbf{k},\omega) &=& U n + \frac{U^2 n(1-n)}{\omega-(1-n)\varepsilon(\mathbf{k})} .
\end{eqnarray}



Long formulas in appendices must be numbered separately, and the formulas
number must include a capital roman letter identifying
the appendix, for examples (A1), (A2), (A3), etc.

\subsection{Figures}

Figure is placed nearby the text place, where it is first cited.
All figures must be cited in consecutive order in the manuscript text.
Since all articles are printed  in the black color, so if colored illustrations or/and figures are submitted, be sure that they will be well resolved in the printed version. However, the colored illustrations or/and figures can be displayed with full colors in the online version, and they are indicated by (Color online) in the figure caption. The figure captions should be descriptive.
Layout figures should size reasonable.

Figures can be prepared in any figure format that is compatible with the LaTex graphics package, for instance pdf, pgn, eps etc. They can directly be embedded within the manuscript text, for example
\begin{figure}[h]
 	\centering
 	\includegraphics[width=0.5\textwidth]{fig1.pdf}
 	\caption{\label{fig1} Figure caption,
 		solid line is theory data, $\Diamond$ is experimental data.}
\end{figure}

If the manuscript is prepared by Microsoft Word, use Times New Roman, size 10 pt for figures captions. The figures captions begin with {\bf Fig. figures number} followed by period, for example, {\bf Fig. 1.}, {\bf Fig. 2.} ... The captions must end with period.

\subsection{References}

Each reference in References should be numbered according to the order of appearance in the main text.
The references numbers are written by Arabic within square brackets. References in languages other than English must be referred to by an English translation with adding to the end a note of original language. References must end with period.


\paragraph*{Reference to a journal publication}
Reference to a published paper must include the names of all the authors, the title of the paper in Italic, the abbreviated journal name, the volume in boldface, the year of publication in parentheses, the first pages number or the article number and the \textit{DOI link or url if available}. When the number of authors is more than 5, the reference can include only the first author, followed by {\it et al.} in Italic.
For example~\cite{streda1989quantised, molenkamp1990quantum, abudinen2021precise},
\bibexample{1}{P. Streda, {\it Quantised thermopower of a channel in the ballistic regime},
\href{https://doi.org/10.1088/0953-8984/1/5/021}{{\it J. Phys. Condens. Matter} \textbf{1} (1989) 1025}.}
\bibexample{2}{L. W. Molenkamp, H. van Houten, C. W. J. Beenakker, R. Eppenga, and C. T. Foxon, {\it Quantum oscillations in thetransverse voltage of a channel in the nonlinear transport regime},
\href{https://doi.org/https://doi.org/10.1103/PhysRevLett.127.211801}{{\it Phys. Rev. Lett.} \textbf{65} (1990) 1052.}
}
\bibexample{3}{F. Abudinen {\it et al.}, {\it Precise Measurement of the $D^0$ and $D^+$ Lifetimes at Belle II},
\href{https://doi.org/https://doi.org/10.1103/PhysRevLett.127.211801}{{\it Phys. Rev. Lett.} \textbf{127} (2021) 211801.}}

\paragraph*{Reference to a book}
Reference to a book must include the names of all the authors, the title, the name of the publisher, the place of the publisher and the year of publication, for example~\cite{grabert2013single},
\bibexample{4}{H. Grabert and M. H. Devoret, \textit{Single charge tunneling: Coulomb blockade phenomena in nanostructures}, vol. 294. 	Springer Science \& Business Media, 2013.}

\paragraph*{Reference to a book in a series}

Reference to a book must include the names of all the authors, the title, the name of a series, the name of the publisher, the place of the publisher and the year of publication, for example~\cite{Santal'o1976},
\bibexample{5}{L. A. Santal\'o, \textit{Integral geometry and geometric probability}, in \textit{Encyclopedia of Mathematics and its Applications},
G. C. Rota, ed., vol. 1, Addison-Wesley, (1976).}

\paragraph*{Reference to a conference/workshop proceeding}

Reference to conference/workshop proceedings must include the names of all the authors, the name of the conference/workshop, the name of the publisher, and the first page of the cited work, the year of publication in parentheses and the \textit{DOI link or url if available}. When the number of authors is more than 5, the reference can include only the first author, for example~\cite{thompson1998magnetic,Pham2011}:
\bibexample{6}{P. Thompson, O. Gutfleisch, J. N. Chapman, I. R. Harris and W. Nicholson, \textit{A magnetic and compositional study of
the disproportionated stage of the solid-hddr process in NdFeB-type materials}, vol. 281, pp. 12–16, Elsevier, 1998, \href{https://doi.org/https://doi.org/10.1016/S0925-8388(98)00762-2}{DOI}.}
\bibexample{7}{M.-T. Pham \textit{et al.}, \textit{A new distance for scale-invariant 3D shape recognition and registration}, in \textit{ICCV}, 2011, \href{https://doi.org/https://doi.org/10.1109/ICCV.2011.6126236}{DOI}.}

\paragraph*{Reference to preprints}
Reference to a preprint must include the names of all the authors, the preprint identifer, the
the year of publication in parentheses and hyperlink for the DOI, for example~\cite{Walter},
\bibexample{8}{W. D. Goldberger and M. B. Wise, \href{https://doi.org/10.48550/arXiv.hep-ph/9907447}{arXiv:9907447} (1999).}

\paragraph*{Reference to websites}
Reference to a website must include the website name, the website http address, the accessed date in parentheses, for example~\cite{CancerResearchUK}
\bibexample{9}{Cancer Research UK, {\it Cancer statistics for the UK},\newline \href{https://www.cancerresearchuk.org/health-professional/cancer-statistics-for-the-uk (accessed 20 October 2021).}{https://www.cancerresearchuk.org/health-professional/cancer-statistics-for-the-uk (accessed 20 October 2021)}.}



\subsection{Footnotes}

Footnotes should be cited where appropriate and
should be cited in consecutive numerical order. They are placed at the bottom of the page on
which they appear. One can use footnotes LaTex command \verb|\footnote|.

If the manuscript is prepared by Microsoft Word, use Times New Roman, size 9 point, single spaced for footnotes.


\section{Conclusions}

The main body of the manuscript ends with Conclusions. The Conclusions summarize the most important results of the research works. Authors are requested not to include reference citation in the Conclusions.


\section*{Acknowledgements}
Acknowledgements are written in a section, which is not numbered and placed after Conclusions, and before the statement of Authors contributions, the statement of Conflict of interest, and References. It is suggested to be one paragraph. Acknowledgment of financial support is listed at the end of Acknowledgements.

\section*{Authors contributions}
When necessary, authors can declare the authors contributions. An example for the statement of the authors contributions as follows

Author 1: Methodology, Investigation, Funding acquisition. Author 2: Formal analysis. Author 3: Formal analysis, Supervision. Author 4: Formal analysis. Casen Panaitescu: Supervision. etc.

\section*{Conflict of interest}
When necessary, authors can declare the statement of conflict of interest, for example:\\
The authors have no conflict of interest to declare.\\
The authors declare that they have no competing financial interests.



\bibliographystyle{cip-v3-bst-submit}
\bibliography{cip-v3-bst-references}

%=========================================================================
\section*{Appendix}

Appendices, if needed, appear after References.


\end{document}

